\uuid{y1Lp}
\exo7id{4136}
\titre{exo7 4136}
\auteur{quercia}
\organisation{exo7}
\datecreate{2010-03-11}
\isIndication{false}
\isCorrection{false}
\chapitre{Fonction de plusieurs variables}
\sousChapitre{Dérivée partielle}
\cpge{mpsi-mp2i.fonctions-deux-variables.b, pcsi.fonctions-deux-variables.b, ptsi.fonctions-deux-variables.b, mpsi-mp2i.fonctions-deux-variables.c, pcsi.fonctions-deux-variables.c, ptsi.fonctions-deux-variables.c, bcpst1.derivees-primitives-integrales.b, bcpst1.fonctions-deux-variables.c}
\module{Analyse}
\niveau{L2}
\difficulte{}

\contenu{
\texte{
Calculer les dérivées partielles des fonctions :
}
\begin{enumerate}
    \item \question{$f(x,y,z) = (x+z)^{(y^z)}$\quad}
    \item \question{$f(x,y) = \min(x,y^2)$  \qquad}
    \item \question{$f(x,y) = \begin{cases}\frac {g(x)-g(y)}{x-y} &\text{ si } x \ne y \cr
                     g'(x) &\text{ si } x = y.\end{cases}$}
\end{enumerate}
}
